\section{Finite covers and domain restrictions}
\label{sec:finite-covers}

The continuation of a determinant depends on the permitted paths.
A germ may have finitely many branches under paths in one domain even
when paths that leave the domain produce infinitely many additional values.
We determine the least cover for the first choice and the components
created by the second.  The definitions of \(\Delta_m,c_m,\omega_m,\alpha_m\)
and their local charts are those of Section~\ref{sec:prime-continuation},
including \(m=1\).  Write \(\mathcal G_m=\operatorname{im}\chi_m\).
A finite branched cover of degree \(d\) is a proper surjective holomorphic
map of connected Riemann surfaces with discrete branch values and total
local degree \(d\) over every point.  Proper means that inverse images
of compact sets are compact; the total local degree is the sum of the
ramification indices in a fibre.

\subsection{An exact criterion for a finite cover}

Let \(D\subset\Hp\) be a connected open set, and put
\(D^\circ=D\setminus S_m\).
Choose \(b\in D^\circ\) and a path in \(\Omega_m\) from \(2\) to \(b\).
It fixes a nonzero continued germ \(f_b\).
Restrict the differential and define
\[
 \mathcal G_m(D)=\left\{\exp\int_\gamma\omega_m:
       \gamma\text{ is a loop in }D^\circ\text{ based at }b\right\}.
\]
Changing the initial path changes \(f_b\) by a constant in \(\mathcal G_m\),
but it leaves this subgroup unchanged.
Let \(\mathcal M(D)\) denote the field of single-valued meromorphic
functions on \(D\).
A germ is algebraic over \(\mathcal M(D)\) if it satisfies a nonzero
polynomial with coefficients in that field.
Its algebraic degree is the least degree of such a polynomial.
A polynomial is irreducible if it has no factorization into polynomials
of smaller positive degrees over the same field.

\begin{theorem}[Finite-cover criterion]
\label{mono:thm:criterion}
The following conditions are equivalent.
\begin{enumerate}
\item The group \(\mathcal G_m(D)\) is finite.
\item There is an integer \(k\geq1\) such that
\(k\int_\gamma\omega_m\in2\pi i\mathbb Z\) for every loop in \(D^\circ\).
\item Some positive integer power of \(f_b\) extends meromorphically to \(D\).
\item A connected finite branched cover of \(D\) carries a meromorphic
function agreeing with the continued germ near a point over \(b\).
\item The germ \(f_b\) is algebraic over \(\mathcal M(D)\).
\end{enumerate}
When these conditions hold, let \(k=|\mathcal G_m(D)|\).
Then \(k\) is the least possible power in condition (iii), the least
cover degree in condition (iv), and the algebraic degree in condition (v).
There is \(A_D\in\mathcal M(D)\) such that the minimal cover is the
branched filling of \(y^k=A_D(s)\), with \(y(b)=f_b(b)\).
The polynomial \(X^k-A_D\) is irreducible over \(\mathcal M(D)\).
\end{theorem}

\begin{proof}
A finite subgroup of \(\mathbb C^*\) of size \(k\) equals the group
\(\mu_k\) of \(k\)-th roots of unity.
Indeed, multiplying all its elements before and after multiplication
by one element \(g\) gives \(g^k=1\).
There are exactly \(k\) roots, so the inclusion is equality.
Conversely a subgroup annihilated by \(k\) has at most \(k\) elements.
This proves the equivalence of (i) and (ii).

Under (ii), define on \(D^\circ\)
\[
 A_D(s)=f_b(b)^k\exp\left(k\int_b^s\omega_m\right).
\]
The period condition makes it path independent and holomorphic.
At a puncture \(a\in D\cap S_m\), a small positive loop gives
\(k\alpha_m(a)\in\mathbb Z\).
Theorem~\ref{mono:prop:local} consequently makes \(A_D\) meromorphic there.
This proves (iii).  Conversely, continuing a meromorphic power around a
loop gives \(g^k=1\), so (iii) implies (ii).

For \(k=|\mathcal G_m(D)|\), the equation \(y^k=A_D\) over \(D^\circ\)
is a connected \(k\)-sheeted cover.
Local roots give its sheets, and continuation realizes every element
of \(\mu_k\), so every sheet is reached from the specified one.
For completeness, at a zero or pole write
\(A_D(s)=(s-a)^v u(s)\) and choose \(h^k=u\).
If \(d=\gcd(k,|v|)\), its local charts are
\[
 s=a+t^{k/d},\qquad
 y=\xi\,t^{v/d}h(a+t^{k/d}),\qquad \xi^k=1.
\]
There are \(d\) distinct charts: changing \(t\) by a \((k/d)\)-th
root identifies exactly \(k/d\) choices of \(\xi\).
They fill all punctured sheets and give a meromorphic \(y\).
When \(v=0\), use \(d=k\) and ordinary unramified charts.
These local changes glue to the given cover on \(D^\circ\).
Above a sufficiently small closed disk there are finitely many compact
coordinate disks.  A finite covering of a compact subset of \(D\)
by such disks proves properness.  This proves (iv).
The polynomial equation also proves (v).

Suppose (iv) holds on a cover \(\pi:X\to D\) of degree \(d\),
with meromorphic function \(F\) and possibly additional branch values
in \(D^\circ\).
Move the base point slightly to avoid them.
Every loop in \(D^\circ\) can be perturbed to avoid the discrete extra
branch values without changing its period.
Lifting such a loop sends the initial sheet to one of at most \(d\) sheets.
The continued value is multiplied by its monodromy factor.
Distinct factors give distinct nonzero values at that point.
Thus \(|\mathcal G_m(D)|\leq d\), proving (i) and the lower bound for cover degree.
The identity \(dF/F=\pi^*\omega_m\), initially valid near the chosen
point, extends as a meromorphic identity on the connected cover.
Thus \(F\) has no zero or pole above \(D^\circ\), and its agreement
with the continued germ persists along the lifted paths.

Finally suppose a polynomial of degree \(d\) vanishes on the germ.
Choose a point in \(D^\circ\) where its leading coefficient is nonzero
and all coefficients are finite.  Such points exist because the leading
coefficient is not identically zero and its exceptional sets are discrete.
Analytic continuation of the polynomial identity shows that all monodromy
values at that point are its roots.
Again these values are distinct, so \(|\mathcal G_m(D)|\leq d\).
This proves (v) implies (i), and every annihilating polynomial has degree
at least \(k\).
The polynomial \(X^k-A_D\) attains that degree and is irreducible:
otherwise one of its smaller-degree factors would vanish on the germ.
A power is single-valued exactly when it annihilates \(\mu_k\), which
requires divisibility by \(k\).  All three minimality claims follow.
\end{proof}

The criterion uses periods of loops in \(D\), not only residues at
points inside \(D\).
For example, choose a small disk about \(a=1/p\) containing no other
point of \(S_m\), where \(p\geq m\geq2\) is prime.
An annulus \(r<|s-a|<R\) inside this disk has no divisor points.
Nevertheless its central loop has residue integral \(2\pi i/p\),
so its minimal cover has degree \(p\).
Deleting the center has removed the local pole from the domain but
has retained the loop around it.

\subsection{Order one, higher orders, and finite restrictions}

For \(m=1\), divisor cancellation gives
\(c_1(1)=1\) and \(c_1(n)=0\) for \(n>1\).
Consequently
\[
 \Delta_1=\zeta^{-1},\qquad
 \omega_1=-\frac{\zeta'}{\zeta}\,ds,\qquad
 \alpha_1(a)=-\ord_a\zeta.
\]
The inverse zeta function is meromorphic on \(\Hp\), and the
scalar continuation of Theorem~\ref{thm:tate} extends it meromorphically
to \(\mathbb C\).
Every residue is an integer and \(\mathcal G_1(D)=\{1\}\) for every connected
\(D\subset\Hp\).  All three degrees in Theorem~\ref{mono:thm:criterion}
are one.  In particular, \(\alpha_1(1/p)=0\) for every prime \(p\).
At \(1\) it has a simple zero, and at a zeta zero of multiplicity \(v\)
it has a pole of order \(v\).

\begin{corollary}[Algebraic obstruction for higher orders]
\label{mono:thm:global}
For \(m\geq2\), the determinant germ is transcendental over
\(\mathcal M(\Hp)\).  No connected finite branched cover of \(\Hp\)
carries its coherent continuation as a meromorphic function, and no
nonzero single-valued meromorphic multiplier changes this conclusion.
\end{corollary}

\begin{proof}
Theorem~\ref{thm:infinite-degree} gives monodromy of order \(p\) at
\(1/p\) for every prime \(p\geq m\).  The image is therefore infinite,
and Theorem~\ref{mono:thm:criterion} excludes both algebraicity and a
finite branched realization.  A meromorphic multiplier adds integer
local orders and integral multiples of \(2\pi i\) to the periods,
so its continuation character is trivial.
\end{proof}

For a reduced rational number \(x\), let \(\operatorname{den}(x)\) be
its positive denominator, with \(\operatorname{den}(0)=1\).
For \(\epsilon>0\), define the finite integer
\[
 M_m(\epsilon)=\operatorname{lcm}
 \{\operatorname{den}(c_m(n)):n\geq1,\ n\epsilon<1\},
\]
where the least common multiple of an empty set is one.

\begin{theorem}[Exact degree on a smaller half-plane]
\label{mono:thm:halfplane}
For continuation along paths in \(\Hp_\epsilon\setminus S_m\),
\[
 \mathcal G_m(\Hp_\epsilon)=\mu_{M_m(\epsilon)}.
\]
Hence the least cover degree and algebraic degree on \(\Hp_\epsilon\)
are exactly \(M_m(\epsilon)\).
Every relatively compact connected domain \(D\subset\Hp\) admits a
finite branched realization of its chosen continued germ.
Its degree divides \(M_m(\epsilon)\) whenever \(D\subset\Hp_\epsilon\).
\end{theorem}

\begin{proof}
Choose \(N\geq m\) with \(1/(N+1)\leq\epsilon\).
For every \(n\leq N\) satisfying \(n\epsilon\geq1\),
\(\zeta(ns)=e^{Z(ns)}\) is holomorphic and nonzero on \(\Hp_\epsilon\).
Its logarithm can be absorbed into the holomorphic remainder.
The only logarithmic factors left have \(n\epsilon<1\).
Every residue is therefore an integer combination of the corresponding
\(c_m(n)\).
The residue-period formula on this convex half-plane puts the monodromy
inside \(\mu_{M_m(\epsilon)}\).
Conversely, the point \(1/n\) lies inside precisely when \(n\epsilon<1\).
Its loop gives \(e^{2\pi i c_m(n)}\).
The subgroup generated by these finitely many roots has order equal to
the least common multiple of their orders: an integer annihilates all
of them precisely when every reduced denominator divides it.
This proves equality.

A compact subset of \(\Hp\) has positive minimum real part.
Thus the closure of a relatively compact \(D\) lies in some
\(\Hp_\epsilon\).  Each loop in \(D\setminus S_m\) is also a loop
in that half-plane, so \(\mathcal G_m(D)\) is a subgroup of the finite cyclic group.
Its order divides \(M_m(\epsilon)\).
Apply Theorem~\ref{mono:thm:criterion}.
\end{proof}

For \(m=2\), the coefficients begin
\[
 c_2(2)=\tfrac12,\quad c_2(3)=\tfrac13,\quad c_2(4)=0,
 \quad c_2(5)=\tfrac15,\quad c_2(6)=-\tfrac16.
\]
Thus the exact half-plane degrees are
\[
 \begin{array}{c|ccccc}
 \epsilon&1/3&1/4&1/5&1/6&1/7\\\hline
 M_2(\epsilon)&2&6&6&30&30
 \end{array}
\]
The strict inequality excludes \(n=3\) at \(\epsilon=1/3\).
The zero coefficient at four contributes no ramification.
At \(1/6\), the exponent \(-1/6\) gives a simple pole after the
sixfold local substitution.
These examples exhibit boundary exclusion, cancellation, and a pole
within the same determinant.

\subsection{Restricting every globally continued branch}

Let \(D\subset\Hp\) be connected, with the path and base point chosen
as in Theorem~\ref{mono:thm:criterion}.
The preimage \(\pi^{-1}(D^\circ)\subset\mathcal E_m\) includes every
branch obtained by paths in the full \(\Omega_m\).
It can have more than one component.  The component containing the
chosen point uses only paths in \(D^\circ\).

\begin{theorem}[Components after restriction]
\label{mono:thm:restriction}
The connected components of \(\pi^{-1}(D^\circ)\) are indexed by the
cosets \(\mathcal G_m/\mathcal G_m(D)\).
Each component is a copy of the minimal pointed continuation cover for
paths in \(D^\circ\), with degree \(|\mathcal G_m(D)|\).
On a component indexed by a representative \(g\in \mathcal G_m\), the determinant
is the chosen component's function multiplied by \(g\).
These identifications extend over the local branched fillings in \(D\).

If \(m\geq2\) and \(D\) is relatively compact, there are infinitely many
components, each of finite degree.
A finite cover carrying a single meromorphic function cannot encode all
globally obtained branch values over \(D\).
Nevertheless the minimal finite cover of \(D\) supports, as separate
functions, all the multiples \(g f\), \(g\in \mathcal G_m\).
\end{theorem}

\begin{proof}
At \(b\), identify the fiber with \(\mathcal G_m\) using the chosen initial path.
Lifting loops in \(D^\circ\) moves a point exactly through its coset
of \(\mathcal G_m(D)\).
Points in the same coset are connected by these lifts.
A path in the preimage connecting two fiber points projects to a loop
in \(D^\circ\), so points in different cosets cannot be connected.
Every component meets that fiber by lifting a path to \(b\).
This proves the component classification.

Prepending a loop with character \(g\) defines a holomorphic map between
the chosen component and its translate.
It has the inverse obtained from \(g^{-1}\), commutes with the projection,
and multiplies the determinant by \(g\).
This proves the stated identification of functions and covers.
At a puncture their exponents are the same \(\alpha_m(a)\).
Multiplication by \(g\) changes only the nonzero constant in
Theorem~\ref{mono:prop:local}, so the maps extend across its local charts.
A local cycle stays within its component, because its loop belongs to
\(\mathcal G_m(D)\).  Filling cannot merge different components.

For relatively compact \(D\), Theorem~\ref{mono:thm:halfplane} makes
\(\mathcal G_m(D)\) finite, whereas Corollary~\ref{mono:thm:global} makes \(\mathcal G_m\) infinite.
The coset set is therefore infinite.
At any ordinary point there are infinitely many distinct global branch
values, all nonzero scalar multiples indexed by \(\mathcal G_m\).
One meromorphic function on a degree-\(d\) cover has at most \(d\)
values above an ordinary point.  It cannot contain that entire set.
On the other hand, each separately specified function \(g f\) is
meromorphic on the same finite cover as \(f\).
This proves both assertions without identifying the two meanings of
carrying all branches.
\end{proof}

For example, on a small disk about \(2\), the group \(\mathcal G_m(D)\) is
trivial and a single local function suffices for paths confined to the disk.
For \(m\geq2\), the global surface restricts to infinitely many disjoint
copies of that disk, with the functions \(g\Delta_m\), \(g\in \mathcal G_m\).
Going around a distant reciprocal-prime point and returning to the disk
changes which copy contains the endpoint.

Theorem~\ref{thm:infinite-degree} also excludes a finite-degree cover
of a neighborhood of zero carrying the specified germs at all sufficiently
small positive reciprocal-prime points.  That conclusion uses their
unbounded local indices.  It gives no natural-boundary statement at any
other point of the imaginary axis.

\subsection{Compatibility with the finite root equations}

The finite degree can also be read directly from the equations
\(y^{M_N}=A_{m,N}(s)\) in Section~\ref{sec:prime-continuation}.
The comparison depends on the initial branch as well as the domain.

\begin{proposition}[The pointed component on a finite half-plane]
\label{prop:finite-root-comparison}
Fix \(m\geq2\), \(N\geq m\), and \(D=U_N=\Hp_{1/(N+1)}\).
Choose \(b=2\) and the initial value \(f_b(2)=\Delta_m(2)\).
Then
\[
 M_m(1/(N+1))=M_N,\qquad A_D=A_{m,N}.
\]
The minimal cover of Theorem~\ref{mono:thm:criterion} identifies with the
finite root surface by \((s,y)\mapsto(s,y)\), including its branched filling.
Inside the global surface, this is the component containing the initial
point.  The components over \(U_N\) are indexed by
\(\mathcal G_m/\mu_{M_N}\).
Changing the initial germ by a factor \(g\in\mathcal G_m\) changes
\(A_D\) by \(g^{M_N}\).
\end{proposition}

\begin{proof}
The strict condition \(n/(N+1)<1\) is exactly \(n\leq N\).
Coefficients below \(m\) vanish, so the two least common multiples agree.
The finite product is the explicitly defined meromorphic function
\[
 A_{m,N}(s)=\exp(M_NT_{m,N}(s))
       \prod_{n=m}^N\zeta(ns)^{-M_Nc_m(n)}.
\]
Both \(A_D\) and \(A_{m,N}\) have logarithmic derivative \(M_N\omega_m\)
on \(U_N\setminus S_m\).
At \(2\), both equal \(\Delta_m(2)^{M_N}\), by the initial extraction
formula.  Their ratio has derivative zero on this connected domain and
value one at \(2\), proving equality there and meromorphically on \(U_N\).
The root equations and initial roots coincide, and their local charts use
the same reduced residue.  The identity map therefore extends over all
filled points.
Theorem~\ref{mono:thm:halfplane} identifies the subgroup with \(\mu_{M_N}\),
and Theorem~\ref{mono:thm:restriction} gives the component set.
Finally the definition of \(A_D\) contains the factor \(f_b(2)^{M_N}\).
Replacing that value by \(g f_b(2)\) multiplies \(A_D\) by \(g^{M_N}\).
Thus paths leaving \(U_N\) can select another component, while paths
confined to \(U_N\) retain the pointed one.
\end{proof}

